\scanpage{182}
% archive foliation: 181
Pratiquement, cela revient à contracter $i_0(k_0),i_1(k_1),i_2(k_2)$
en des pts, et à se borner à des cocycles continus sur le
\deleted{\unsure{complexe}}\inserted{complexe simplicial}
\uline{ponctué} obtenu
\diagram{\begin{tikzcd}
\overline K_0&
\overline K_1\arrow[l,shift left=.5ex]\arrow[l,shift right=.5ex]&
\overline K_2\arrow[l,shift left=1ex]\arrow[l]\arrow[l,shift right=1ex]
\end{tikzcd}}

On utilise le fait que
\[
\pi_0(k_*)=e
\]
\[
``\pi_1(k_*)=e''\quad\text{i.e. }\Pi^1(k_*,G)=e
\]
donc \uline{tte donnée de descente admet une ponctuation}
et on utilisera le fait
\[
\pi_0(K_*)=e
\]
pour le fait que \deleted{\unsure{deux ponctuations}}
tout \uline{automorphisme} d'une structure de descente
\uline{ponctuée} est trivial.
\deleted{[\unsure{et est même d'une}
\unsure{ponctuation donnée, que deux ponctuations, ne diffèrent que par un
élémt de $G$, d'ailleurs unique}]}

\deleted{Désignons \unsure{les}}
